=o(t^{-L}).
\end{equation}
\end{lemma}

\begin{proof}
Lemma~\ref{lem:local-multidim-quotient-taylor-remainder} gives the
coefficient bound and the \(|z|\le1\) part of
\eqref{eq:multidim-uniform-taylor-remainder-bound}.

For the large-\(z\) bound, fix \(|z|>1\) and split according to whether
\(|y-z|\le |z|/2\).  If \(|y-z|\le |z|/2\), then \(|y|\le3|z|/2\), and
\[
F_y(z)\le (1+|y|^2)^a\le C_d|z|^{d+1}.
\]
On the complementary region \(|y-z|>|z|/2\), the elementary inequality
\[
1+|y|^2
\le 2(1+|y-z|^2+|z|^2)
\le 2(1+|z|^2)(1+|y-z|^2)
\]
gives
\[
F_y(z)\le 2^a(1+|z|^2)^a\le C_d|z|^{d+1}.
\]
Thus
\[
\sup_y F_y(z)\le C_d|z|^{d+1}\le C_{d,L}|z|^L
\qquad(|z|>1),
\]
where \(L\ge d+1\) is used only in the final comparison.  The boundedness of
the Taylor coefficients gives, for \(|z|>1\),
\[
\sup_y\sum_{|\alpha|\le L}|U_\alpha^{(d)}(y)|\,|z|^{|\alpha|}
\le C_{d,L}\sum_{k=0}^L |z|^k
\le C_{d,L}'|z|^L.
\]
Consequently
\[
\sup_y\left(
|F_y(z)|
+\sum_{|\alpha|\le L}|U_\alpha^{(d)}(y)|\,|z|^{|\alpha|}
\right)
\le C_{d,L}|z|^L
\qquad(|z|>1),
\]
and this proves the \(|z|>1\) part of
\eqref{eq:multidim-uniform-taylor-remainder-bound}.

With \(z=X/t\), the bound just proved gives, uniformly in \(y\),
\[
|r_t(y)|
\le
C_{d,L}\left[
t^{-(L+1)}\EE\!\left(|X|^{L+1}\mathbf1_{\{|X|\le t\}}\right)
+t^{-L}\EE\!\left(|X|^L\mathbf1_{\{|X|>t\}}\right)
\right].
\]
The two terms have the required order because
\[
t^{-(L+1)}\EE\!\left(|X|^{L+1}\mathbf1_{\{|X|\le t\}}\right)
=
t^{-L}\EE\!\left[
|X|^L\frac{|X|}{t}\mathbf1_{\{|X|\le t\}}
\right]
=o(t^{-L})
\]
by dominated convergence, and
\[
t^{-L}\EE\!\left(|X|^L\mathbf1_{\{|X|>t\}}\right)=o(t^{-L})
\]
by integrability of \(|X|^L\).  These estimates are uniform in \(y\), so they
prove the \(L^\infty\)-part of
\eqref{eq:multidim-uniform-density-remainder-consequence}.  The
\(L^1(\Omega_d)\)-part follows from the same uniform bound because
\(\Omega_d\) is a probability measure.
\end{proof}

\medskip
\noindent\textbf{Multidimensional radial Poisson jet hypothesis
\((\mathrm{MPJ}_{d,L})\).}
For \(d\ge1\) and an integer \(L\), say that \((\nu,d,L)\) satisfies
\((\mathrm{MPJ}_{d,L})\) if
\[
L\ge\max\{2,d+1\},\qquad
\bar{\gamma}=\int_{\RR^d}\gamma\,\dd\nu(\gamma)\ \text{exists},\qquad
\EE|\gamma-\bar{\gamma}|^L<\infty .
\]
The condition \(L\ge d+1\) is part of the analytic hypothesis, not a
notational convention: it is exactly what makes the large-translation bound
in \eqref{eq:multidim-uniform-taylor-remainder-bound} integrable after
averaging over \(X_c/t\).
\medskip

\begin{lemma}[Where the multidimensional finite-\(L\) hypothesis enters]%
\label{lem:multidim-threshold-proof-interface}
In the proof of
Theorems~\ref{thm:multidimensional-density-quotient} and
\ref{thm:multidimensional-poisson-entropy-jets}, the restriction
\(L\ge d+1\) is used only to close the global
\(L^\infty_y\) quotient-remainder estimate in
Lemma~\ref{lem:multidim-uniform-radial-poisson-taylor}.  More precisely, the
large-translation part of the deterministic Taylor remainder is bounded by
\[
\sup_{y\in\RR^d}
\left|
F_y(z)-\sum_{|\alpha|\le L}U_\alpha^{(d)}(y)z^\alpha
\right|
\le C_{d,L}|z|^L
\qquad(|z|>1),
\]
only because \(F_y(z)\) itself can have size \(C_d|z|^{d+1}\) near the
translated pole \(y=z\).  After setting \(z=X_c/t\), the proof therefore
requires
\[
t^{-L}\EE\!\left[
|X_c|^L\mathbf1_{\{|X_c|>t\}}
+|X_c|^L\frac{|X_c|}{t}\mathbf1_{\{|X_c|\le t\}}
\right]=o(t^{-L}),
\]
which is exactly the finite \(L\)th absolute-moment input used in
\eqref{eq:multidim-uniform-density-remainder-consequence}.

Thus the multidimensional entropy statements below are packaged as
finite-\(L\) jet theorems.  The coefficient calculation in
Proposition~\ref{prop:multidim-formal-covariance-layer} is separate from this
remainder step: it is a finite-covariance identity for the formal quadratic
layer.
When \((\mathrm{MPJ}_{d,L})\) is also available, that identity becomes the
leading \(t^{-4}\) coefficient of the finite multidimensional entropy jet by
Corollary~\ref{cor:multidim-leading-kl-coefficient}.
\end{lemma}

\begin{proposition}[Large-translation obstruction for the quotient method]%
\label{prop:multidim-large-translation-obstruction}
Fix \(d\ge1\), put \(a=(d+1)/2\), and let
\[
F_y(z)=
\left(\frac{1+|y|^2}{1+|y-z|^2}\right)^a .
\]
For an integer \(M\ge0\), set
\[
T_M(y,z):=\sum_{|\alpha|\le M}U_\alpha^{(d)}(y)z^\alpha,
\qquad
U_\alpha^{(d)}(y)=\frac1{\alpha!}\partial_z^\alpha F_y(0).
\]
There are constants \(c_{d,M},C_{d,M}>0\) and \(R_{d,M}<\infty\) such that,
for all \(|z|\ge R_{d,M}\),
\begin{equation}\label{eq:multidim-obstruction-linfty-lower}
\|F_\cdot(z)-T_M(\cdot,z)\|_{L^\infty(\RR^d)}
\ge c_{d,M}|z|^{d+1},
\end{equation}
and
\begin{equation}\label{eq:multidim-obstruction-l2-lower}
\|F_\cdot(z)-T_M(\cdot,z)\|_{L^2(\Omega_d)}
\ge c_{d,M}|z|^{(d+1)/2}.
\end{equation}
Consequently the \(L^\infty\)-quotient proof cannot be lowered below the
\(d+1\) tail exponent by a sharper bookkeeping of the same Taylor remainder.
Replacing the uniform norm by an \(L^2(\Omega_d)\) norm improves the
deterministic large-translation exponent only to \((d+1)/2\).  This leaves a
possible low-dimensional refinement in \(d\le3\), but it does not give an
all-dimensional finite-covariance quotient theorem, since in dimensions
\(d\ge4\) the required deterministic exponent is still stronger than
covariance.  Thus Theorems~\ref{thm:multidimensional-density-quotient} and
\ref{thm:multidimensional-poisson-entropy-jets} remain full finite-\(L\)
quotient-jet theorems under \((\mathrm{MPJ}_{d,L})\).  The covariance
calculation in Proposition~\ref{prop:multidim-formal-covariance-layer} does not
use this obstruction; it identifies the formal quadratic covariance layer
itself under the finite second-moment assumption.
\end{proposition}

\begin{proof}
It is enough to work on the ball
\[
B_z:=\{y\in\RR^d:\ |y-z|\le 1/4\}.
\]
For \(y\in B_z\) and \(|z|\ge2\), one has
\(|z|/2\le |y|\le 2|z|\) and \(1+|y-z|^2\le 17/16\).  Hence
\[
F_y(z)
\ge c_d(1+|y|^2)^a
\ge c_d'|z|^{d+1}.
\]
On the same ball the Taylor coefficients at the origin are much smaller than
the translated peak.  Indeed, the derivative structure
\eqref{eq:multidim-poisson-derivative-structure}, with \(w=0\), gives
\[
|U_\alpha^{(d)}(y)|
\le C_{d,\alpha}(1+|y|)^{-|\alpha|}
\qquad(|y|\ge1),
\]
and therefore
\[
|U_\alpha^{(d)}(y)z^\alpha|
\le C_{d,\alpha}
\qquad (y\in B_z,\ |z|\ge2).
\]
For fixed \(M\), \(|T_M(y,z)|\le C_{d,M}\) on \(B_z\).  Increasing
\(R_{d,M}\) if necessary, the translated peak dominates this bounded Taylor
polynomial, and
\[
|F_y(z)-T_M(y,z)|\ge c_{d,M}|z|^{d+1}
\qquad (y\in B_z,\ |z|\ge R_{d,M}).
\]
This proves the \(L^\infty\) lower bound.

For the \(L^2(\Omega_d)\) bound, use the same pointwise lower bound on
\(B_z\) and the radial Poisson weight
\(\Omega_d(\dd y)=c_d(1+|y|^2)^{-(d+1)/2}\dd y\).  Since
\(|z|/2\le |y|\le2|z|\) on \(B_z\),
\[
\Omega_d(B_z)\ge c_d''|z|^{-(d+1)}.
\]
Thus
\[
\|F_\cdot(z)-T_M(\cdot,z)\|_{L^2(\Omega_d)}^2
\ge c_{d,M}|z|^{2(d+1)}\Omega_d(B_z)
\ge c_{d,M}'|z|^{d+1},
\]
which is \eqref{eq:multidim-obstruction-l2-lower} after taking square roots.
\end{proof}

\begin{theorem}[Finite-\(L\) multidimensional density quotient expansion, \(L\ge d+1\)]%
\label{thm:multidimensional-density-quotient}
Let \(d\ge1\), let
\[
P_t^{(d)}(x):=c_d\,\frac{t}{(t^2+|x|^2)^{(d+1)/2}},
\qquad x\in\RR^d,\ t>0,
\]
be the radial Poisson kernel on \(\RR^d\), normalized to integrate to one.
Let \(\nu\) be a probability measure on \(\RR^d\).  Fix an integer \(L\)
such that \((\nu,d,L)\) satisfies
\((\mathrm{MPJ}_{d,L})\), and put
\(X_c:=\gamma-\bar{\gamma}\).  Set
\[
h_t:=P_t^{(d)}*\nu,\qquad
g_t(x):=P_t^{(d)}(x-\bar{\gamma}).
\]
For multiindices \(\alpha\in\NN_0^d\), write
\[
\mu_\alpha:=\EE[X_c^\alpha]\quad(|\alpha|\le L).
\]
Let
\[
\Omega_d(\dd y):=P_1^{(d)}(y)\dd y
\]
and define the universal functions \(U_\alpha^{(d)}\) by the Taylor expansion
at \(z=0\)
\begin{equation}\label{eq:multidim-poisson-ratio-taylor}
\left(\frac{1+|y|^2}{1+|y-z|^2}\right)^{(d+1)/2}
=\sum_{\alpha\in\NN_0^d}U_\alpha^{(d)}(y)z^\alpha .
\end{equation}
Then, for \(x=\bar{\gamma}+ty\), the Cauchy-scale density quotient satisfies
\begin{equation}\label{eq:multidim-density-expansion}
\frac{h_t(\bar{\gamma}+ty)}{g_t(\bar{\gamma}+ty)}-1
=\sum_{j=2}^{L}D_j^{(d)}(y;\nu)t^{-j}
+\rho_{L,t}(y),
\qquad
\|\rho_{L,t}\|_\infty+\|\rho_{L,t}\|_{L^1(\Omega_d)}=o(t^{-L}),
\end{equation}
where, for \(2\le j\le L\),
\begin{equation}\label{eq:multidim-density-layers}
D_j^{(d)}(y;\nu):=
\sum_{|\alpha|=j}\mu_\alpha U_\alpha^{(d)}(y).
\end{equation}
Lemma~\ref{lem:multidim-uniform-radial-poisson-taylor} supplies, uniformly
in \(y\), both the finite Taylor remainder bound
\eqref{eq:multidim-uniform-taylor-remainder-bound} and its
\(L^\infty\cap L^1(\Omega_d)\) integrated consequence
\eqref{eq:multidim-uniform-density-remainder-consequence}; these are the only
large-\(z\) and density-remainder estimates used here.
\end{theorem}

\begin{proof}
Write \(a=(d+1)/2\).  For \(x=\bar{\gamma}+ty\),
\[
\frac{P_t^{(d)}(x-\gamma)}{g_t(x)}
=
\left(\frac{t^2+|x-\bar{\gamma}|^2}{t^2+|x-\gamma|^2}\right)^a
=
\left(\frac{1+|y|^2}{1+|y-X_c/t|^2}\right)^a.
\]
Thus
\[
\widetilde\delta_t(y)
:=\frac{h_t(\bar{\gamma}+ty)}{g_t(\bar{\gamma}+ty)}-1
=\EE\!\left[
\left(\frac{1+|y|^2}{1+|y-X_c/t|^2}\right)^a-1
\right].
\]
Apply Lemma~\ref{lem:multidim-uniform-radial-poisson-taylor} with
\(X=X_c\).  Its uniform Taylor-remainder consequence gives the error term
below in both
\(L^\infty(\RR^d)\) and \(L^1(\Omega_d)\).  Since \(U_0^{(d)}=1\) and the
first-order terms integrate to zero against the centred variable \(X_c\), we
obtain \eqref{eq:multidim-density-expansion}.  The proof of the stated
remainder is exactly the tail split in the Taylor lemma, with \(z=X_c/t\):
the region \(|X_c|\le t\) gives the integrable factor
\(|X_c|^L(|X_c|/t)\), while \(|X_c|>t\) gives the finite \(L\)th-moment tail.
\end{proof}

\begin{lemma}[Standalone multidimensional density-to-entropy transfer]%
\label{lem:multidim-density-to-entropy-transfer}
Let \((S,\Omega)\) be a probability space, let \(L\ge2\), and let
\(D_2,\ldots,D_L\in L^\infty(\Omega)\).  Suppose that, for all sufficiently
large \(t\),
\[
\delta_t=\sum_{j=2}^{L}D_jt^{-j}+\rho_t,
\]
where
\[
\|\rho_t\|_\infty+\|\rho_t\|_{L^1(\Omega)}=o(t^{-L})
\]
and
\[
\|\delta_t\|_\infty=O(t^{-2}).
\]
Then, for
\[
\Phi(s)=(1+s)\log(1+s)-s,
\]
\begin{equation}\label{eq:abstract-multidim-transfer-coefficients}
\mathfrak A_m(D_2,\ldots,D_L)
:=
\sum_{r=2}^{\lfloor m/2\rfloor}
\frac{(-1)^r}{r(r-1)}
\sum_{\substack{2\le j_1,\ldots,j_r\le L\\
j_1+\cdots+j_r=m}}
\int_S D_{j_1}\cdots D_{j_r}\,\dd\Omega
\end{equation}
one has
\begin{equation}\label{eq:abstract-multidim-transfer-expansion}
\int_S\Phi(\delta_t)\,\dd\Omega
=
\sum_{m=4}^{L+2}\mathfrak A_m(D_2,\ldots,D_L)t^{-m}
+o(t^{-(L+2)}).
\end{equation}
\end{lemma}

\begin{proof}
Put
\[
B_t:=\sum_{j=2}^{L}D_jt^{-j}.
\]
Since the \(D_j\)'s are bounded and \(\Omega\) has total mass one,
\(\|B_t\|_\infty=O(t^{-2})\), and every finite product appearing in
\eqref{eq:abstract-multidim-transfer-coefficients} belongs to
\(L^1(\Omega)\).  The assumed smallness of \(\delta_t\) gives
\(|\delta_t|\le1/2\) for large \(t\).  Let
\[
R:=\left\lfloor\frac{L+2}{2}\right\rfloor .
\]
For \(|s|\le1/2\),
\[
\Phi(s)=\sum_{r=2}^{R}\frac{(-1)^r}{r(r-1)}s^r+O(|s|^{R+1}).
\]
Because \(\|\delta_t\|_\infty=O(t^{-2})\), the Taylor tail has
\[
\int_SO(|\delta_t|^{R+1})\,\dd\Omega
=O(t^{-2(R+1)})=o(t^{-(L+2)}).
\]

It remains to replace \(\delta_t^r\) by \(B_t^r\), for \(2\le r\le R\).  In
\[
\delta_t^r-B_t^r=(B_t+\rho_t)^r-B_t^r
\]
each summand contains at least one factor \(\rho_t\).  Hence, using one
\(\rho_t\) in \(L^1(\Omega)\) and all other factors in \(L^\infty\),
\[
\|\delta_t^r-B_t^r\|_{L^1(\Omega)}
\le
\sum_{\ell=1}^{r}\binom r\ell
\|B_t\|_\infty^{r-\ell}
\|\rho_t\|_\infty^{\ell-1}
\|\rho_t\|_{L^1(\Omega)}
=o(t^{-L}t^{-2(r-1)}).
\]
Since \(r\ge2\), this is \(o(t^{-(L+2)})\).  Therefore
\[
\int_S\Phi(\delta_t)\,\dd\Omega
=
\sum_{r=2}^{R}\frac{(-1)^r}{r(r-1)}
\int_SB_t^r\,\dd\Omega
+o(t^{-(L+2)}).
\]

Finally expand the finite polynomial:
\[
\int_SB_t^r\,\dd\Omega
=
\sum_{m=2r}^{Lr}t^{-m}
\sum_{\substack{2\le j_1,\ldots,j_r\le L\\
j_1+\cdots+j_r=m}}
\int_SD_{j_1}\cdots D_{j_r}\,\dd\Omega .
\]
The terms with \(m\le L+2\), summed over \(2\le r\le R\), give exactly the
coefficients in \eqref{eq:abstract-multidim-transfer-coefficients}.  The
remaining finite terms have integer exponents \(m\ge L+3\), hence contribute
\(O(t^{-(L+3)})=o(t^{-(L+2)})\).  This proves
\eqref{eq:abstract-multidim-transfer-expansion}.
\end{proof}

\begin{theorem}[Finite-\(L\) multidimensional entropy transfer, \(L\ge d+1\)]%
\label{thm:multidimensional-poisson-entropy-jets}
Assume the hypotheses and notation of
Proposition~\ref{thm:multidimensional-density-quotient}.  Thus
\((\nu,d,L)\) satisfies \((\mathrm{MPJ}_{d,L})\), \(h_t=P_t^{(d)}*\nu\),
\(g_t(x)=P_t^{(d)}(x-\bar{\gamma})\), and the density layers
\(D_j^{(d)}\) are those in \eqref{eq:multidim-density-layers}.  For
\(4\le m\le L+2\), define \(\mathcal A_m^{(d)}(\nu)\) by
\begin{equation}\label{eq:multidim-entropy-coefficient-layers}
\mathcal A_m^{(d)}(\nu)=
\sum_{r=2}^{\lfloor m/2\rfloor}
\frac{(-1)^r}{r(r-1)}
\sum_{\substack{2\le j_1,\ldots,j_r\le L\\
j_1+\cdots+j_r=m}}
\int_{\RR^d}D_{j_1}^{(d)}(y;\nu)\cdots
D_{j_r}^{(d)}(y;\nu)\,\Omega_d(\dd y).
\end{equation}
Then
\begin{equation}\label{eq:multidim-poisson-entropy-jet}
D_{\rm KL}(h_t\|g_t)
=\sum_{m=4}^{L+2}\mathcal A_m^{(d)}(\nu)t^{-m}
+o(t^{-(L+2)})
\qquad(t\to\infty).
\end{equation}
\end{theorem}

\begin{proof}
The change of variables gives
\[
g_t(\bar{\gamma}+ty)t^d\dd y=P_1^{(d)}(y)\dd y=\Omega_d(\dd y),
\]
and therefore, with
\(\widetilde\delta_t(y)=h_t(\bar{\gamma}+ty)/g_t(\bar{\gamma}+ty)-1\),
the density masses cancel in the same variables:
\[
\int_{\RR^d}\widetilde\delta_t(y)\,\Omega_d(\dd y)
=\int_{\RR^d}\bigl(h_t(x)-g_t(x)\bigr)\,\dd x=0.
\]
Hence
\[
D_{\rm KL}(h_t\|g_t)
=\int_{\RR^d}\Phi(\widetilde\delta_t(y))\,\Omega_d(\dd y),
\qquad
\Phi(s)=(1+s)\log(1+s)-s.
\]
By Proposition~\ref{thm:multidimensional-density-quotient},
\[
\widetilde\delta_t(y)
=\sum_{j=2}^{L}D_j^{(d)}(y;\nu)t^{-j}
+\rho_{L,t}(y),
\qquad
\|\rho_{L,t}\|_\infty+\|\rho_{L,t}\|_{L^1(\Omega_d)}=o(t^{-L}).
\]
The finite sum starts at order \(t^{-2}\).  For each \(2\le j\le L\), the
retained layer \(D_j^{(d)}\) is bounded because it is a finite linear
combination of the bounded Taylor modes \(U_\alpha^{(d)}\) from
Lemma~\ref{lem:multidim-uniform-radial-poisson-taylor}, with finite
coefficients \(\mu_\alpha\), \(|\alpha|=j\).  Together with
\(\|\rho_{L,t}\|_\infty=o(t^{-L})\), this gives
\(\|\widetilde\delta_t\|_\infty=O(t^{-2})\).  Therefore the tuple
\[
(S,\Omega,\delta_t,D_2,\ldots,D_L,\rho_t)
=
(\RR^d,\Omega_d,\widetilde\delta_t,
D_2^{(d)}(\cdot;\nu),\ldots,D_L^{(d)}(\cdot;\nu),\rho_{L,t})
\]
satisfies every hypothesis of the standalone transfer lemma
Lemma~\ref{lem:multidim-density-to-entropy-transfer}.  The abstract
coefficients
\(\mathfrak A_m(D_2^{(d)},\ldots,D_L^{(d)})\) in
\eqref{eq:abstract-multidim-transfer-coefficients} are precisely the
\(\mathcal A_m^{(d)}(\nu)\) defined in
\eqref{eq:multidim-entropy-coefficient-layers}.  The lemma therefore proves
\eqref{eq:multidim-poisson-entropy-jet}.
\end{proof}

\begin{corollary}[Multidimensional coefficient formula]%
\label{cor:multidim-coefficient-formula}
Assume the hypotheses and notation of
Proposition~\ref{thm:multidimensional-density-quotient}.  For \(4\le m\le L+2\),
the coefficient in Theorem~\ref{thm:multidimensional-poisson-entropy-jets}
is equivalently
\begin{equation}\label{eq:multidim-entropy-coefficient}
\mathcal A_m^{(d)}(\nu):=
\sum_{r=2}^{\lfloor m/2\rfloor}
\frac{(-1)^r}{r(r-1)}
\sum_{\substack{|\alpha_1|,\ldots,|\alpha_r|\ge2\\
|\alpha_1|+\cdots+|\alpha_r|=m\\
|\alpha_j|\le L}}
\mu_{\alpha_1}\cdots\mu_{\alpha_r}
\int_{\RR^d}\prod_{j=1}^rU_{\alpha_j}^{(d)}(y)\,\Omega_d(\dd y).
\end{equation}
The coefficient is finite and universal: it depends on \(d\), \(m\), and the
centred tensor moments
\(\{\mu_\alpha:2\le|\alpha|\le m-2\}\), but not on any higher moment of
\(\nu\).  In particular a contribution of total entropy degree \(m\) cannot
contain a tensor moment of order above \(m-2\).
\end{corollary}

\begin{proof}
Substitute the density-layer definition
\eqref{eq:multidim-density-layers} into
\eqref{eq:multidim-entropy-coefficient-layers} and collect multiindices by
homogeneous degree.  Since every entropy monomial has at least two density
factors, each individual multiindex in a total degree \(m\) contribution has
length at most \(m-2\).
\end{proof}

\begin{lemma}[Multidimensional parity and low-order display]%
\label{lem:multidim-parity-low-order-display}
Assume the hypotheses and notation of
Theorems~\ref{thm:multidimensional-density-quotient} and
\ref{thm:multidimensional-poisson-entropy-jets}.  Then all odd entropy
coefficients vanish:
\begin{equation}\label{eq:multidim-odd-parity-cancellation}
\mathcal A_{2q+1}^{(d)}(\nu)=0
\qquad(2q+1\le L+2).
\end{equation}
The first potentially nonzero coefficients are
\begin{align}
\mathcal A_4^{(d)}
&=\frac12\int_{\RR^d}D_2^2\,\dd\Omega_d,
\label{eq:multidim-first-coeff-a4}\\
\mathcal A_6^{(d)}
&=\int_{\RR^d}D_2D_4\,\dd\Omega_d
+\frac12\int_{\RR^d}D_3^2\,\dd\Omega_d
-\frac16\int_{\RR^d}D_2^3\,\dd\Omega_d,
\label{eq:multidim-first-coeff-a6}\\
\mathcal A_8^{(d)}
&=\int_{\RR^d}D_2D_6\,\dd\Omega_d
+\int_{\RR^d}D_3D_5\,\dd\Omega_d
+\frac12\int_{\RR^d}D_4^2\,\dd\Omega_d
\notag\\
&\quad
-\frac12\int_{\RR^d}D_2^2D_4\,\dd\Omega_d
-\frac12\int_{\RR^d}D_2D_3^2\,\dd\Omega_d
+\frac1{12}\int_{\RR^d}D_2^4\,\dd\Omega_d,
\label{eq:multidim-first-coeff-a8}
\end{align}
with the convention that a displayed term containing \(D_j\) is omitted when
\(j>L\).  Here \(D_j=D_j^{(d)}(\,\cdot\,;\nu)\).

The leading term is coercive in the covariance:
Proposition~\ref{prop:multidim-formal-covariance-layer} gives an explicit positive
decomposition of \(\mathcal A_4^{(d)}\), so
\(\mathcal A_4^{(d)}\ge0\) and equality holds exactly for zero covariance.
No fixed sign is asserted for the higher even coefficients; the formulas
\eqref{eq:multidim-first-coeff-a6}--\eqref{eq:multidim-first-coeff-a8}
display the alternating Bregman-Taylor contributions that destroy a general
positivity statement beyond the leading quadratic layer.

For even \(m\), the tensor-moment order \(m-2\) enters the universal
coefficient \(\mathcal A_m^{(d)}\) only through the quadratic-layer pairing
\(\int D_2D_{m-2}\,\dd\Omega_d\).  This is the algebraic sharpness statement
for the displayed universal coefficient; the detected part of the top tensor
is the component paired with \(D_2\).
\end{lemma}

\begin{remark}[Notation for the leading multidimensional layer]
The finite-jet coefficient in
\eqref{eq:multidim-first-coeff-a4} is denoted
\(\mathcal A_4^{(d)}\) because it is part of the full
\((\mathrm{MPJ}_{d,L})\) entropy jet.  The same quadratic expression is
denoted \(\mathcal Q_d\) in
Proposition~\ref{prop:multidim-formal-covariance-layer} when only finite
covariance is assumed.  Corollary~\ref{cor:multidim-leading-kl-coefficient}
is the separate step that turns this formal layer into a leading KL coefficient
under \((\mathrm{MPJ}_{d,L})\).
\end{remark}

\begin{proof}[Proof of the parity and low-order display]
The identity \(F_{-y}(z)=F_y(-z)\) in
\eqref{eq:multidim-poisson-ratio-taylor} implies
\[
U_\alpha^{(d)}(-y)=(-1)^{|\alpha|}U_\alpha^{(d)}(y),
\qquad
D_j^{(d)}(-y;\nu)=(-1)^jD_j^{(d)}(y;\nu).
\]
Since \(\Omega_d\) is even, every product in
\eqref{eq:multidim-entropy-coefficient-layers} with odd total degree has zero
integral.  This proves \eqref{eq:multidim-odd-parity-cancellation}.  The
formulas \eqref{eq:multidim-first-coeff-a4}--\eqref{eq:multidim-first-coeff-a8}
are obtained by listing the ordered compositions of \(m\) by parts at least
two and applying the Taylor coefficients \(1/2,-1/6,1/12\) of \(\Phi\).
For even \(m\), the only compositions containing a part \(m-2\) are
\((2,m-2)\) and \((m-2,2)\); their combined contribution is
\(\int D_2D_{m-2}\,\dd\Omega_d\).  All remaining compositions use only
density layers of order at most \(m-3\).
\end{proof}

\begin{proof}[Moment identities used in Proposition~\ref{prop:multidim-formal-covariance-layer}]
Let \(Y\) have density \(\Omega_d(\dd y)=P_1^{(d)}(y)\dd y\), set
\(R=|Y|^2\), and put
\[
D_d:=(d+1)(d+3)(d+5)(d+7).
\]
Then
\[
\EE\frac{1}{(1+R)^4}=\frac{105}{D_d},\qquad
\EE\frac{R}{(1+R)^4}=\frac{15d}{D_d},\qquad
\EE\frac{R^2}{(1+R)^4}=\frac{3d(d+2)}{D_d}.
\]
Also, if \(\Theta\) is uniformly distributed on the unit sphere in \(\RR^d\)
and \(A\) is a real symmetric traceless \(d\times d\) matrix, then
\[
\EE_\Theta\bigl[(\Theta^{\mathsf T}A\Theta)^2\bigr]
=\frac{2}{d(d+2)}\|A\|_{\rm HS}^2 .
\]
Indeed, since \(\Omega_d\) is radial, the polar decomposition
\(Y=\sqrt R\,\Theta\) has \(\Theta\) uniform on the sphere, independent of
\(R\), and \(R\sim{\rm BetaPrime}(d/2,1/2)\).  Hence
\begin{equation}\label{eq:radial-beta-prime-moments}
\EE\frac{R^k}{(1+R)^4}
=
\frac{B\bigl(d/2+k,\,9/2-k\bigr)}
{B\bigl(d/2,\,1/2\bigr)}
\qquad(k=0,1,2).
\end{equation}
Expanding the beta functions gives the three displayed rational moments; in
particular the \(d=1,\ k=0\) value is \(35/128\).
The angular identity follows by contracting
\[
\EE_\Theta[\Theta_i\Theta_j\Theta_k\Theta_l]
=\frac{\delta_{ij}\delta_{kl}
+\delta_{ik}\delta_{jl}
+\delta_{il}\delta_{jk}}{d(d+2)}
\]
with the entries of \(A\), and then using \(\operatorname{tr}A=0\).
\end{proof}


\begin{proposition}[Algebraic covariance quadratic form under finite covariance]%
\label{prop:multidim-formal-covariance-layer}
Let \(d\ge1\), and let \(\nu\) be a probability measure on \(\RR^d\) with
finite centred second moment.  Put
\[
\bar{\gamma}=\int_{\RR^d}\gamma\,\dd\nu(\gamma),
\qquad
X_c=\gamma-\bar{\gamma},
\]
and let
\[
\Sigma:=\EE[X_cX_c^{\mathsf T}]
\]
be the centred covariance matrix.  Put
\[
\Sigma_0:=\Sigma-\frac{\operatorname{tr}\Sigma}{d}I_d .
\]
Let \(b_\Sigma\) denote the degree-two centred density mode obtained from the
radial Poisson quotient, explicitly
\begin{equation}\label{eq:multidim-leading-theorem-density-mode}
b_\Sigma(y)
 =
\frac{d+1}{2}\left(
-\frac{\operatorname{tr}\Sigma}{1+|y|^2}
+\frac{(d+3)y^{\mathsf T}\Sigma y}{(1+|y|^2)^2}
\right).
\end{equation}
The algebraic quadratic entropy-integrand layer determined by this mode is
the finite quantity
\begin{equation}\label{eq:multidim-leading-positive-defect}
\mathcal Q_d(\Sigma)
:=\frac12\int_{\RR^d}b_\Sigma(y)^2\,\Omega_d(\dd y)
=
c_d^{\rm iso}\,(\operatorname{tr}\Sigma)^2
+c_d^{\rm tr}\,\|\Sigma_0\|_{\rm HS}^2,
\end{equation}
where
\begin{equation}\label{eq:multidim-leading-positive-constants}
c_d^{\rm iso}:=
\frac{3(d+1)(7d+9)}
{4d(d+3)(d+5)(d+7)},
\qquad
c_d^{\rm tr}:=
\frac{3(d+1)(d+3)}
{4(d+5)(d+7)} .
\end{equation}
Two immediate checks are
\[
d=1:\qquad \Sigma_0=0,\qquad c_1^{\rm iso}=\frac18,\qquad
\mathcal Q_1(\Sigma)=\frac{(\operatorname{tr}\Sigma)^2}{8},
\]
and, for an isotropic covariance \(\Sigma=\tau I_d\),
\[
\Sigma_0=0,\qquad
\mathcal Q_d(\Sigma)
=c_d^{\rm iso}d^2\tau^2 .
\]
This identity assumes only \(\EE|X_c|^2<\infty\) and is a covariance-layer
calculation.  Finite covariance alone identifies the quadratic form below;
the present proof of the KL asymptotic uses the stronger quotient-jet
hypothesis \((\mathrm{MPJ}_{d,L})\), supplied here by finite centred
\(L\)-moment with \(L\ge d+1\).  The asymptotic statement
\[
D_{\rm KL}\bigl(P_t^{(d)}*\nu\,\|\,P_t^{(d)}(\cdot-\bar{\gamma})\bigr)
=\mathcal Q_d(\Sigma)t^{-4}+o(t^{-4})
\]
is Corollary~\ref{cor:multidim-leading-kl-coefficient}.
